\section{שיעור 08 — משוואות דיפרנציאליות חלקיות: סדר ראשון וסיווג סדר שני}

\begin{summarybox}{נושאי השיעור}
בשיעור זה עסקנו בפתרון משוואות דיפרנציאליות חלקיות (מד"ח) מסדר ראשון, מעבר למשוואות קוואזי-ליניאריות, והקדמה למשוואות מסדר שני (משוואת הגלים).
\begin{itemize}
    \item שיטת הקרקטריסטיקות למד"ח מסדר ראשון (ליניאריות וקוואזי-ליניאריות).
    \item בעיית קושי (Cauchy Problem) ותנאים לקיום פתרון יחיד.
    \item פתרון משוואות לא הומוגניות על ידי העלאת מימד (Implicit Method).
    \item מבוא למשוואת הגלים החד-ממדית.
    \item סיווג בעיות שפה: דיריכלה, נוימן וקושי.
\end{itemize}
\end{summarybox}

\subsection{מד"ח מסדר ראשון ושיטת הקרקטריסטיקות}

בשיעור הקודם הצגנו את הצורה הכללית של משוואה ליניארית הומוגנית מסדר ראשון בשני משתנים:
\[
a(x,y)u_x + b(x,y)u_y = 0
\]
המשמעות הגיאומטרית היא שהנגזרת הכיוונית של $u$ בכיוון הווקטור $\vec{v}=(a,b)$ מתאפסת. מכאן נובע כי הפתרון $u(x,y)$ קבוע לאורך עקומות המשיקות לשדה הווקטורי הזה. עקומות אלו נקראות \textbf{קרקטריסטיקות} \LR{(Characteristics)}.

\begin{definitionbox}{משוואת הקרקטריסטיקות}
עבור השדה המוגדר על ידי המקדמים $(a,b)$, משוואת הקרקטריסטיקות נתונה על ידי המד"ר:
\[
\frac{dx}{a(x,y)} = \frac{dy}{b(x,y)}
\]
פתרון המשוואה נותן משפחת עקומות $\xi(x,y) = C$. הפתרון הכללי של המד"ח הוא פונקציה כללית של הקרקטריסטיקות:
\[
u(x,y) = f(\xi(x,y))
\]
\end{definitionbox}

\subsubsection{בעיית קושי (Cauchy Problem)}
בעיית קושי עבור מד"ח מסדר ראשון דורשת למצוא פתרון $u(x,y)$ המקיים את המשוואה ועובר דרך עקומה נתונה $\Gamma$ במרחב $(x,y,u)$.
פורמלית, נתון הערך של $u$ על עקומה במישור $xy$:
\[
u(x,y)\big|_{\Gamma} = \phi(s)
\]
\begin{theorembox}{תנאי לקיום וייחוד פתרון}
כדי שבעיית קושי תהיה מוגדרת היטב ותניב פתרון יחיד בסביבת העקומה, נדרש שהעקומה $\Gamma$ עליה נתונים תנאי ההתחלה \textbf{תחתוך} את הקרקטריסטיקות (ולא תשיק להן או תתלכד איתן). כלומר, העקומה צריכה לחצות כל קרקטריסטיקה פעם אחת בלבד באזור הרלוונטי.
\end{theorembox}

\begin{examplebox}{דוגמה 1: משוואה ליניארית עם מקדמים קבועים}
נתונה המשוואה:
\[ u_x + 2u_y = 0 \]
עם תנאי שפה על ציר ה-$y$: $u(0,y) = y^2$.

\textbf{פתרון:}
\begin{enumerate}
    \item \textbf{מציאת הקרקטריסטיקות:}
    \[ \frac{dx}{1} = \frac{dy}{2} \implies 2dx = dy \implies y - 2x = C \]
    הקרקטריסטיקות הן קווים ישרים בשיפוע 2.
    \item \textbf{פתרון כללי:}
    \[ u(x,y) = f(y - 2x) \]
    \item \textbf{הצבת תנאי שפה:}
    נציב $x=0$:
    \[ u(0,y) = f(y) = y^2 \]
    קיבלנו את הצורה הפונקציונלית של $f$. כעת נציב את הארגומנט המלא $y-2x$:
    \[ u(x,y) = (y - 2x)^2 \]
\end{enumerate}
\end{examplebox}

\begin{examplebox}{דוגמה 2: פתרון גיאומטרי}
נתונה המשוואה:
\[ x u_x - y u_y = 0 \]
משוואת הקרקטריסטיקות:
\[ \frac{dx}{x} = \frac{dy}{-y} \implies \ln|x| = -\ln|y| + C \implies xy = \text{const} \]
הקרקטריסטיקות הן היפרבולות. הפתרון הכללי הוא $u(x,y) = f(xy)$.
אם נתון תנאי על מעגל היחידה $x^2+y^2=1$, יש לוודא שהמעגל חותך את ההיפרבולות בצורה טרנסברסלית (Transversal).
\end{examplebox}

\subsection{משוואות קוואזי-ליניאריות (Quasi-Linear PDEs)}

משוואה קוואזי-ליניארית היא משוואה בה הנגזרות מופיעות בצורה ליניארית, אך המקדמים יכולים להיות תלויים בפונקציה הנעלמת $u$ עצמה. הצורה הכללית ב-$n$ משתנים:
\[
\sum_{i=1}^n a_i(x_1, \dots, x_n, u) \frac{\partial u}{\partial x_i} = c(x_1, \dots, x_n, u)
\]
במקרה כזה, לא ניתן למצוא את הקרקטריסטיקות במישור $x$ בלבד, כיוון שהן תלויות בפתרון $u$.

\subsubsection{שיטת הפתרון הסתום (Implicit Method)}
הרעיון הוא לחפש פתרון בצורה סתומה \LR{(implicit form)}:
\[ w(x_1, \dots, x_n, u) = 0 \]
נשתמש בזהויות הנגזרת של פונקציה סתומה:
\[ \frac{\partial u}{\partial x_i} = - \frac{\partial w / \partial x_i}{\partial w / \partial u} \]
בהצבה למשוואה המקורית וסידור מחדש, מקבלים משוואה הומוגנית ליניארית במימד $n+1$:
\[
\sum_{i=1}^n a_i \frac{\partial w}{\partial x_i} + c \frac{\partial w}{\partial u} = 0
\]
כעת ניתן לכתוב מערכת משוואות קרקטריסטיות מורחבת:

\begin{resultbox}{מערכת הקרקטריסטיקות למשוואה קוואזי-ליניארית}
\[
\frac{dx_1}{a_1} = \frac{dx_2}{a_2} = \dots = \frac{dx_n}{a_n} = \frac{du}{c}
\]
\end{resultbox}
מערכת זו מגדירה עקומות במרחב ה-$n+1$ ממדי $(x_1, \dots, x_n, u)$. הפתרון הכללי מתקבל על ידי מציאת $n$ אינטגרלים ראשוניים בלתי תלויים $\phi_k(x,u) = C_k$, והפתרון הכללי הוא קשר פונקציונלי ביניהם:
\[ \Phi(\phi_1, \dots, \phi_n) = 0 \]

\begin{examplebox}{דוגמה: משוואה קוואזי-ליניארית}
פתור את המשוואה:
\[ x u_x + y u_y = u \]
עם תנאי שפה על הפרבולה $y=x^2$: $u(x, x^2) = \sin x$.

\textbf{פתרון:}
\begin{enumerate}
    \item \textbf{משוואות קרקטריסטיות:}
    \[ \frac{dx}{x} = \frac{dy}{y} = \frac{du}{u} \]
    \item \textbf{מציאת אינטגרלים ראשוניים:}
    מתוך $\frac{dx}{x} = \frac{dy}{y}$ נקבל $\ln x = \ln y + C \implies \frac{y}{x} = c_1$.
    מתוך $\frac{dx}{x} = \frac{du}{u}$ נקבל $\frac{u}{x} = c_2$.
    \item \textbf{פתרון כללי:}
    הפתרון הסתום הוא $\Phi(\frac{y}{x}, \frac{u}{x}) = 0$, או בצורה מפורשת:
    \[ \frac{u}{x} = g\left(\frac{y}{x}\right) \implies u = x g\left(\frac{y}{x}\right) \]
    \item \textbf{הצבת תנאי קושי:}
    על העקומה $y=x^2$, נתון $u=\sin x$. נציב בפתרון הכללי:
    \[ \sin x = x g\left(\frac{x^2}{x}\right) = x g(x) \]
    מכאן אנו מוצאים את הפונקציה $g$:
    \[ g(x) = \frac{\sin x}{x} \]
    \item \textbf{פתרון סופי:}
    נציב חזרה את הארגומנט $y/x$ לתוך $g$:
    \[ u(x,y) = x \cdot \frac{\sin(y/x)}{y/x} = \frac{x^2}{y} \sin\left(\frac{y}{x}\right) \]
\end{enumerate}
\end{examplebox}

\subsection{מבוא למשוואות מסדר שני: משוואת הגלים}

המוטיבציה העיקרית ללימוד מד"ח בפיזיקה היא משוואות מסדר שני, ובראשן משוואת הגלים החד-ממדית:
\[
u_{tt} - c^2 u_{xx} = 0
\]
זוהי משוואה היפרבולית. ניתן לפרק אותה לאופרטורים מסדר ראשון:
\[
\left(\frac{\partial}{\partial t} - c\frac{\partial}{\partial x}\right)\left(\frac{\partial}{\partial t} + c\frac{\partial}{\partial x}\right)u = 0
\]
פירוק זה מרמז על קיום שתי משפחות של קרקטריסטיקות:
\[
x - ct = \text{const}, \quad x + ct = \text{const}
\]
הפתרון הכללי של משוואת הגלים (פתרון ד'אלמבר) הוא סכום של גל נע ימינה וגל נע שמאלה:
\[
u(x,t) = f(x - ct) + g(x + ct)
\]

\subsubsection{תנאי שפה והתחלה למשוואת הגלים}
בניגוד למד"ח מסדר ראשון, עבור משוואת הגלים נדרש לספק מידע רב יותר כדי לקבל פתרון יחיד.
עבור בעיית קושי ב-$t=0$ (תנאי התחלה), יש לספק:
\begin{enumerate}
    \item את צורת הגל: $u(x,0)$.
    \item את המהירות הרגעית של המיתר: $u_t(x,0)$.
\end{enumerate}
מבחינה פיזיקלית, נתונים אלו מקבילים למיקום ותנע במכניקה קלאסית.

\begin{notebox}{תחום תלות (Domain of Dependence)}
המידע בנקודה $(x,t)$ תלוי אך ורק בתנאי ההתחלה בתוך חרוט האור (התחום התחום על ידי הקרקטריסטיקות היוצאות מ-$(x,t)$ לאחור בזמן). אם נתוני קושי ניתנים על עקומה שאינה חותכת את הקרקטריסטיקות בצורה נכונה (למשל, עקומה המקבילה לקרקטריסטיקה), הבעיה אינה מוגדרת היטב.
\end{notebox}

\subsection{סיווג משוואות ליניאריות מסדר שני}

הצורה הכללית ביותר של מד"ח ליניארית מסדר שני בשני משתנים היא:
\[
A u_{xx} + 2B u_{xy} + C u_{yy} + D u_x + E u_y + F u = G
\]
סיווג המשוואה נקבע על ידי החלק הראשי (הנגזרות השניות), בדומה לסיווג חתכי חרוט, לפי הדיסקרימיננטה $\Delta = B^2 - AC$:

\begin{table}[H]
\centering
\begin{tabular}{@{}llp{6cm}@{}}
\toprule
\textbf{סוג המשוואה} & \textbf{תנאי} & \textbf{דוגמה פיזיקלית} \\ \midrule
היפרבולית \LR{(Hyperbolic)} & $B^2 - AC > 0$ & משוואת הגלים ($u_{tt} - c^2 u_{xx} = 0$) \\
פרבולית \LR{(Parabolic)} & $B^2 - AC = 0$ & משוואת החום/דיפוזיה ($u_t - k u_{xx} = 0$) \\
אליפטית \LR{(Elliptic)} & $B^2 - AC < 0$ & משוואת לפלס ($u_{xx} + u_{yy} = 0$) \\ \bottomrule
\end{tabular}
\caption{סיווג משוואות דיפרנציאליות חלקיות מסדר שני}
\end{table}

\subsubsection{סוגי תנאי שפה}
בהמשך הקורס נדון בהתאמת תנאי השפה לסוג המשוואה:
\begin{itemize}
    \item \textbf{בעיית דיריכלה \LR{(Dirichlet)}:} הפונקציה $u$ נתונה על השפה $\Gamma$.
    \item \textbf{בעיית נוימן \LR{(Neumann)}:} הנגזרת הנורמלית $\frac{\partial u}{\partial n}$ נתונה על השפה. פתרון נקבע עד כדי קבוע.
    \item \textbf{בעיית קושי \LR{(Cauchy)}:} גם $u$ וגם הנגזרת הנורמלית נתונים על השפה (מתאים בעיקר למשוואות היפרבוליות, כמו תנאי התחלה למיתר).
\end{itemize}

\begin{summarybox}{סיכום והשלכות להמשך}
\begin{itemize}
    \item סיימנו את הדיון במשוואות מסדר ראשון. הכלי המרכזי הוא המרת המד"ח למערכת משוואות דיפרנציאליות רגילות (קרקטריסטיקות).
    \item במשוואות קוואזי-ליניאריות, הקרקטריסטיקות תלויות בפתרון עצמו, מה שיכול להוביל לחיתוך קרקטריסטיקות ולהיווצרות "גלי הלם" (Shock waves), אם כי לא העמקנו בכך בקורס זה.
    \item בשיעורים הבאים נתמקד בפתרון משוואות מסדר שני (גלים, חום, לפלס) בתחומים סופיים ואינסופיים, תוך שימוש בשיטות כמו הפרדת משתנים.
\end{itemize}
\end{summarybox}